\documentclass[11pt,a4paper]{article}

\usepackage[utf8]{inputenc}
\usepackage[T1]{fontenc}
\usepackage{amsmath,amssymb,amsthm}
\usepackage[margin=2.0cm]{geometry}
\usepackage{enumitem}
\usepackage{parskip}
\setlength{\parskip}{4pt}
\setlength{\abovedisplayskip}{4pt}
\setlength{\belowdisplayskip}{4pt}
\setlength{\abovedisplayshortskip}{2pt}
\setlength{\belowdisplayshortskip}{2pt}

\newcommand{\R}{\mathbb{R}}
\newcommand{\N}{\mathbb{N}}
\newcommand{\norm}[1]{\lVert #1 \rVert}
\newcommand{\abs}[1]{\lvert #1 \rvert}
\newcommand{\nz}{\norm{\cdot}_2}
\newcommand{\ninf}{\norm{\cdot}_\infty}
\newcommand{\nR}{\norm{\cdot}_{\R}}
\newcommand{\nRz}{\norm{\cdot}_{\R^2}}
\newcommand{\dm}{\,\mathrm{d}}

\title{Einsendeaufgaben Lektion~1 -- Lösungen\\[2pt]
\large 61211 Analysis, WS 2026/2027}
\date{}

\begin{document}
\maketitle
\vspace{-1.2cm}

% =====================================================================
\section*{Aufgabe 1}

\noindent
$\norm{(x_1,x_2)}=\abs{x_1-x_2}+\abs{x_1+x_2}+\abs{x_2}$.

\subsection*{(a) Norm? \hfill (6 P.)}
\textbf{Ja.} Def.~1.4.1 (i)--(iii):
\begin{itemize}[leftmargin=1.6em,itemsep=1pt]
\item \emph{Definitheit:} $\norm{x}\ge0$. Aus $\norm{x}=0$ folgt $\abs{x_2}=0$, also
$x_2=0$, und dann $\abs{x_1}=0$, also $x_1=0$.
\item \emph{Homogenität:} $\norm{\alpha x}=\abs{\alpha x_1-\alpha x_2}+\abs{\alpha x_1+\alpha x_2}+\abs{\alpha x_2}
=\abs{\alpha}\norm{x}$.
\item \emph{Dreiecksungleichung:} Mit den linearen Funktionalen
$\ell_1(x)=x_1-x_2$, $\ell_2(x)=x_1+x_2$, $\ell_3(x)=x_2$ ist
$\norm{x}=\sum_i\abs{\ell_i(x)}$ und
$\abs{\ell_i(x+y)}\le\abs{\ell_i(x)}+\abs{\ell_i(y)}$; Summation liefert
$\norm{x+y}\le\norm{x}+\norm{y}$.
\end{itemize}
Also ist $\norm{\cdot}$ eine Norm.

\subsection*{(b) Konstanten \hfill (4 P.)}
Wegen $\abs{x_1-x_2}+\abs{x_1+x_2}=2\max\{\abs{x_1},\abs{x_2}\}$ gilt mit
$M:=\max\{\abs{x_1},\abs{x_2}\}$:
\[
  \norm{x}=2M+\abs{x_2}.
\]
\emph{Oben:} $M\le\norm{x}_2,\ \abs{x_2}\le\norm{x}_2
\Rightarrow \norm{x}\le2\norm{x}_2+\norm{x}_2=3\norm{x}_2$.

\emph{Unten:} Fall $\abs{x_1}\ge\abs{x_2}$:
$\norm{x}=2\abs{x_1}+\abs{x_2}$ und
\[
  \norm{x}^2-4\norm{x}_2^2=4\abs{x_1}\abs{x_2}-3x_2^2
  =\abs{x_2}\underbrace{(4\abs{x_1}-3\abs{x_2})}_{\ge0}\ge0 .
\]
Fall $\abs{x_2}\ge\abs{x_1}$: $\norm{x}=3\abs{x_2}$ und
$\norm{x}^2-4\norm{x}_2^2=5x_2^2-4x_1^2\ge x_1^2\ge0$.
Also stets $\norm{x}\ge2\norm{x}_2$.

Somit ist $\boxed{c_1=2,\ c_2=3}$ zulässig. Sie sind optimal, da
$\norm{(1,0)}=2\norm{(1,0)}_2$ und $\norm{(0,1)}=3\norm{(0,1)}_2$.

% =====================================================================
\section*{Aufgabe 2}

\noindent
$X=\R^3$, $\norm{x}=\norm{(x_2,x_3)}_{\R^2}+\norm{x_1}_{\R}$, $x=(x_1,x_2,x_3)$.

\subsection*{(a) Normierter Raum \hfill (7 P.)}
$X$ ist ein reeller Vektorraum. Mit $u=(x_2,x_3)$, $v=(y_2,y_3)$:
\begin{itemize}[leftmargin=1.6em,itemsep=1pt]
\item $\norm{x}\ge0$, und
$\norm{x}=0\Leftrightarrow\norm{u}_{\R^2}=0\wedge\norm{x_1}_{\R}=0
\Leftrightarrow u=0\wedge x_1=0\Leftrightarrow x=0$.
\item $\norm{\alpha x}=\norm{\alpha u}_{\R^2}+\norm{\alpha x_1}_{\R}
=\abs{\alpha}\norm{u}_{\R^2}+\abs{\alpha}\norm{x_1}_{\R}=\abs{\alpha}\norm{x}$.
\item $\norm{x+y}=\norm{u+v}_{\R^2}+\norm{x_1+y_1}_{\R}
\le\norm{u}_{\R^2}+\norm{v}_{\R^2}+\norm{x_1}_{\R}+\norm{y_1}_{\R}=\norm{x}+\norm{y}$.
\end{itemize}
Also ist $(X,\norm{\cdot})$ ein normierter Raum.

\subsection*{(b) Äquivalent zu $\nz$? \hfill (3 P.)}
\textbf{Ja.} Da $\R^3$ endlichdimensional ist, sind je zwei Normen äquivalent
(Satz~1.5.11; vgl.\ Hilfssatz~1.5.10). Explizit: Nach Hilfssatz~1.5.10 gibt es
$a,b>0$ mit $a\norm{u}_2\le\norm{u}_{\R^2}\le b\norm{u}_2$; ferner
$\norm{x_1}_{\R}=c\abs{x_1}$ mit $c=\norm{1}_{\R}>0$. Dann
\[
  \norm{x}\le b\norm{u}_2+c\abs{x_1}\le\max\{b,c\}\sqrt2\,\norm{x}_2,
  \qquad
  \norm{x}\ge a\norm{u}_2+c\abs{x_1}\ge\min\{a,c\}\,\norm{x}_2 .
\]

% =====================================================================
\section*{Aufgabe 3}

\noindent
$X=\{f\in C(\R):\lim_{\abs{x}\to\infty}\abs{f(x)}=0\}$.

\subsection*{(a) Normierter Raum \hfill (5 P.)}
\emph{$X$ ist Unterraum von $C(\R)$:} $f+\alpha g$ ist stetig, und
$\abs{f+\alpha g}\le\abs{f}+\abs{\alpha}\abs{g}\to0$ für $\abs{x}\to\infty$.

\emph{Jedes $f\in X$ ist beschränkt:} zu $\varepsilon=1$ gibt es $R>0$ mit
$\abs{f(x)}<1$ für $\abs{x}>R$; auf $[-R,R]$ ist $f$ beschränkt, etwa
$\abs{f}\le M$. Also $\ninf(f)\le\max\{M,1\}<\infty$, d.\,h.\ $X\subseteq B(\R)$.

Nach Beispiel~1.4.3(1) ist $\ninf$ eine Norm auf $B(\R)$; die Einschränkung auf
den Unterraum $X$ ist daher eine Norm auf $X$.

\subsection*{(b) Banachraum? \hfill (5 P.)}
\textbf{Ja.} $(B(\R),\ninf)$ ist nach Beispiel~1.5.19(i) vollständig. Wir zeigen:
$X$ ist abgeschlossen in $B(\R)$. Sei $(f_n)\subseteq X$ mit
$\norm{f_n-f}_\infty\to0$, $f\in B(\R)$.

\emph{$f$ stetig:} zu $x_0,\varepsilon$ wähle $n$ mit
$\norm{f_n-f}_\infty<\tfrac\varepsilon3$ und $\delta$ mit
$\abs{f_n(x)-f_n(x_0)}<\tfrac\varepsilon3$ für $\abs{x-x_0}<\delta$; dann
\[
  \abs{f(x)-f(x_0)}\le\abs{f-f_n}(x)+\abs{f_n(x)-f_n(x_0)}+\abs{f_n(x_0)-f(x_0)}
  <\tfrac\varepsilon3+\tfrac\varepsilon3+\tfrac\varepsilon3=\varepsilon .
\]
\emph{$f$ verschwindet im Unendlichen:} zu $\varepsilon$ wähle $n$ mit
$\norm{f_n-f}_\infty<\tfrac\varepsilon2$ und $R$ mit $\abs{f_n(x)}<\tfrac\varepsilon2$
für $\abs{x}>R$; dann $\abs{f(x)}\le\abs{f-f_n}(x)+\abs{f_n(x)}<\varepsilon$.

Also $f\in X$, d.\,h.\ $X$ abgeschlossen. Ein abgeschlossener Unterraum eines
Banachraums ist vollständig, also ist $(X,\ninf)$ ein Banachraum.

% =====================================================================
\section*{Aufgabe 4}

\noindent
$X=C[a,b]$, $g\in C^1[a,b]$, $g\ge0$, $g'>0$, $\norm{f}_X=\int_a^b g(x)\abs{f(x)}\dm x$.

\subsection*{(a) $x\mapsto\abs{f(x)}$ stetig \hfill (4 P.)}
Für $f\in C[a,b]$ und $x_0\in[a,b]$ gilt
$\bigl\lvert\abs{f(x)}-\abs{f(x_0)}\bigr\rvert\le\abs{f(x)-f(x_0)}$.
Da $f$ in $x_0$ stetig ist, gibt es zu $\varepsilon>0$ ein $\delta>0$ mit
$\abs{f(x)-f(x_0)}<\varepsilon$ für $\abs{x-x_0}<\delta$, also
$\bigl\lvert\abs{f(x)}-\abs{f(x_0)}\bigr\rvert<\varepsilon$. Somit ist
$\abs{f}$ stetig ($\abs{f}=\abs{\cdot}\circ f$).

\subsection*{(b) $\norm{\cdot}_X$ ist eine Norm \hfill (6 P.)}
$g\abs{f}$ ist stetig, also integrierbar (Wohldefiniertheit).

\begin{itemize}[leftmargin=1.6em,itemsep=1pt]
\item \emph{Definitheit:} $\norm{f}_X\ge0$, $\norm{0}_X=0$. Sei
$\norm{f}_X=0$. Dann ist $h=g\abs{f}\ge0$ stetig mit $\int h=0$, also $h\equiv0$
(sonst wäre $h>0$ auf einem Intervall). Wegen $g'>0$ ist $g$ streng monoton
wachsend, also $g>0$ auf $(a,b]$. Damit $f(x)=0$ für $x\in(a,b]$ und mit der
Stetigkeit $f(a)=0$; d.\,h.\ $f\equiv0$.
\item \emph{Homogenität:}
$\norm{\alpha f}_X=\int g\abs{\alpha f}=\abs{\alpha}\int g\abs{f}=\abs{\alpha}\norm{f}_X$.
\item \emph{Dreiecksungleichung:}
$\norm{f+h}_X=\int g\abs{f+h}\le\int g(\abs{f}+\abs{h})
=\int g\abs{f}+\int g\abs{h}=\norm{f}_X+\norm{h}_X$.
\end{itemize}
Damit ist $\norm{\cdot}_X$ eine Norm auf $X$.

\end{document}
